Compactness is a foundational criterion in redistricting law and algorithmic map-drawing,
yet the redistricting literature employs at least six distinct metrics---Polsby-Popper (PP),
Reock, Convex Hull Ratio (CH), Schwartzberg (S), Length-Width Ratio (LW), and
Population-Weighted Compactness (PWC)---each measuring a geometrically distinct
property. This paper provides the first systematic taxonomy of these metrics
applied to the \textsc{bisect} algorithmic redistricting suite, reporting
pairwise Pearson correlations across all congressional districts generated by four
structure algorithms (standard-bisect, prime-factor, ratio-optimal, moving-knife)
for North Carolina ($k=14$, 2{,}195 tracts), Wisconsin ($k=8$, 1{,}406 tracts),
and Texas ($k=38$, 5{,}265 tracts) under 2020 Census data.
We find that area-based metrics (Reock, CH) correlate strongly within-family
($r > 0.75$) while correlating weakly with perimeter-based metrics (PP, S)
across families ($r < 0.5$), and that PWC is largely orthogonal to all
geometric metrics ($r < 0.4$ with PP).
Courts have cited at least four distinct metrics across redistricting cases---PP
in \textit{Rucho v.\ Common Cause} (2019), Reock in North Carolina and Wisconsin
litigation, convex hull in \textit{Shaw v.\ Reno} (1993), and Schwartzberg-equivalent
criteria in Colorado commission proceedings---confirming that practitioners
require a multi-metric profile.
We provide a six-cell decision table mapping each metric to its primary legal
use case and a synthesis of findings from K.1--K.7.
